\section{Arbitrariness of Assignment and Constraints on Acquisition}\label{sec:arb}

\subsection{Two relations that should not be confused}

The claim that the sign is arbitrary has a long history in linguistics. In Saussure's formulation the link between a signifier and what it signifies is not motivated by any intrinsic resemblance \citep{saussure1983}. In Peirce's classification the relation of a sign to its object may be iconic, indexical or symbolic, and the letters of an alphabet are by and large symbols, tied to what they stand for by convention and habit and not by likeness \citep{peirce1931}. For writing, the thesis has a direct and checkable form. There is no necessity that the letter-shape B should serve for the sound it conventionally represents in English, the same language is written in many scripts, and a language's script can change by decree, as the replacement of an Arabic-based script by a Latin-based one in Turkey in 1928 shows \citep{lewis1999}. A small number of scripts retain iconic residues, but in the alphabets of immediate interest the assignment is conventional through and through, and the historical record of the derivation of letter shapes from one another shows convention accumulating on convention \citep{gelb1952,daniels1996}.

What arbitrariness of assignment does not say is that all assignments are equally easy to learn, to tell apart, to typeset, or to read in sequence. Between the set of signs and the reader there is a second structure that the assignment leaves untouched, namely the similarity relation among the shapes themselves. Two letters may be confusable, or share features, or be related by a reflection or a rotation, and this holds whatever sounds or functions are attached to them. The visual substitution alphabet makes the two relations separable. Its assignments are conventional, since there is no reason beyond the experiment's convenience for B to be replaced by P and not by some other letter. But the choice between B and P and the choice between C and K differ in the perceptual bridge they provide, and that difference is a fact about the similarity structure of the shapes and about the reader's habits, not about the assignment as such. The project's theoretical position can therefore be put as follows: the conventionality of writing is demonstrated by the existence of many workable assignments, while the conditions that make a particular convention usable are bodily, perceptual and historical, and a programme of ciphers, typefaces and drifting editions exhibits both claims within a single apparatus.

\subsection{What is known about the acquisition of letter forms}

The literature on how readers acquire letter forms supplies several constraints that bear directly on the design of intermediate alphabets. E. J. Gibson's account of perceptual learning treats the acquisition of an alphabet as the discovery of distinctive features, the invariant properties by which letters differ from one another, and the early experiments with letter-like forms tested which transformations children found easy to discriminate \citep{gibson1969,gibson1962}. On this view the difficulty of learning a new alphabet depends on whether its members differ from one another by features the learner is already sensitive to. A related hypothesis, the cultural recycling of cortical maps, proposes that reading draws on visual machinery that evolved for other purposes and is therefore shaped by prior biases \citep{dehaene2007,dehaene2009}, and there is evidence that the contour configurations of letters across scripts resemble those found in natural scenes \citep{changizi2006}. If these accounts are broadly right, then the set of shapes that can be learned easily is not arbitrary in the way the set of assignments is, and the freedom of a designer is bounded by the reader's perceptual priors.

One constraint is of special relevance to visual ciphers. Children and adults tend to generalise across mirror images, treating a shape and its reflection as the same object, and learning to read involves suppressing this generalisation for the letters where reflection changes identity, such as b and d or p and q \citep{dehaene2010}. A substitution that exchanges mirror-related lower-case letters therefore recruits one of the best documented sources of confusion in letter learning, and a designer who wishes to build resemblance into an alphabet should be aware that resemblance of this kind is not always helpful. It can be a bridge, since the reader's eye is drawn to the right neighbourhood, and it can also be an obstacle, since the reader's eye is drawn to the wrong member of the neighbourhood. Which of these happens for a given pair is an empirical question, and the relative importance of the two effects is one of the quantities the programme in Section~\ref{sec:empirical} is designed to estimate.

\subsection{Levels of knowing a sign}

It is useful to separate several levels at which a reader may know a letter form, because the claims of the project differ by level. At the lowest level a reader recognises a form in context, resolving it through the surrounding word and sentence without being able to say what it is in isolation. At the next level the reader identifies the form when it is presented alone, which requires that its association with a letter be stable and not merely contextually inferred. At a third level the reader can name the form or decode it deliberately, as a student does with a table of signs. At the highest level the reader can produce the form, writing or drawing it correctly from memory. The distinction between receptive and productive knowledge is standard in the study of vocabulary acquisition, where recognition typically runs ahead of production \citep{nation2001}, and the same asymmetry is expected here. A drifting book is designed to exercise the first level and, with luck, the second. It makes no claim to teach the third and the fourth, and any assertion that readers have acquired a script has to say at which level they have acquired it and how this was measured.

\subsection{Precedents}

None of this is without precedent. The Turkish alphabet reform was carried out by a state with the machinery of schooling and legislation behind it, and its theoretical interest for present purposes is that it shows a national convention replaced within a generation \citep{lewis1999}. The Initial Teaching Alphabet used in British schools in the 1960s was a transitional script, intended to be learned first and then replaced by traditional orthography, and the evaluations of that programme are an obvious source of methods and cautions \citep{downing1967}. Languages have long been written in scripts that are not their own: the Arabic-script writing of Romance and of several African languages in the traditions known as Aljamiado and Ajami is a standing demonstration that the same language can ride on different graphic vehicles \citep{daniels1996}. Among speakers of Arabic, informal Latin-script chat orthography assigns digits to Arabic sounds by visual resemblance, a 3 for the letter ayn and a 7 for the letter ha being the commonly cited examples \citep{palfreyman2003}, and this is a naturally occurring visual cipher built by its users for their own purposes. At the level of the individual reader, studies of reading in geometrically transformed typography have found that practice at reading inverted or reflected text leaves traces that survive for a long time, with measurable savings when readers are tested a year later \citep{kolers1976}. All of these show that conventions can be changed and that readers adapt. They differ from the present proposal in one respect that is also the source of its novelty, namely that the acquisition is to take place inside a single volume, without a separate stage of instruction, and the conditions under which this can succeed are not settled by any of the precedents.
